\documentclass[10pt]{article}
\usepackage{tmlr}
\usepackage[T1]{fontenc}
\usepackage[utf8]{inputenc}
\usepackage{amsmath,amssymb,amsthm,booktabs,tabularx,multirow}
\usepackage{graphicx,xcolor,enumitem,placeins,float}
\usepackage{hyperref,url}
\hypersetup{colorlinks=true,linkcolor=blue!45!black,citecolor=blue!45!black,urlcolor=blue!45!black,pdfauthor={},pdftitle={Auditing Sequence Prediction with Known Label Rules}}
\newcommand{\prov}[1]{{\color{black!65}\scriptsize[#1]}}
\newcommand{\orig}{\prov{ORIG}}
\newcommand{\ou}{\prov{ORIG-UNVERIFIED}}
\newcommand{\recomp}{\prov{RECOMP}}
\newcommand{\newres}{\prov{NEW}}
\newcommand{\req}[2]{{\color{brown!75!black}\small[TODO-REQUIRED #1: #2]}}
\newcommand{\opt}[2]{{\color{black!65}\small[TODO-OPTIONAL #1: #2]}}
\newcommand{\ind}{\mathbf{1}}
\newcommand{\AUC}{\mathrm{AUC}}
\newcommand{\R}{\mathbb{R}}
\newtheorem{proposition}{Proposition}
\newtheorem{definition}{Definition}
\setlist[itemize]{leftmargin=*,topsep=2pt,itemsep=1pt,parsep=0pt}
\title{Auditing Sequence Prediction\\with Known Label Rules}
\author{Anonymous authors}
% Anonymous TMLR submission mode; no accepted or preprint option.
\begin{document}
\raggedbottom
\maketitle

\begin{abstract}
A high sequence-prediction score does not show which input features a model uses. We examine what can be established using explicit label rules, restricted feature baselines, and order perturbations. On noisy ordered compliance, lag-pair logistic regression attains AUC $0.6709$ \recomp, compared with $0.6705$ for the true-rule score \newres. Yet we construct opposite-label inputs with identical pair counts, proving that those features cannot represent the clean rule exactly. For hidden-subset parity, a parity-specific learner recovers the subset from $256$ training examples and makes no errors on $100{,}000$ test examples \newres, limiting what neural training failures establish about the task. Key-inversion parity provides an order-sensitive comparison whose label is independent of small fixed sets of positions. Reported fine-tuning and inversion-parity experiments add model-level evidence, subject to verification of their run archives. Clinical diagnosis sequences and music recommendation illustrate the need to preserve visible events when testing order and to distinguish recent transitions from dependence on a long history. Together, the analyses specify which conclusions each comparison supports and which require further evidence. The contribution is a concrete evaluation procedure with explicit limits on its interpretation. All unverified results and unfinished checks remain identified.
\end{abstract}

{\small\textbf{Working-draft status.} Provenance tags distinguish verified original values \orig, unverified reported values \ou, recomputation \recomp, and new analyses \newres. Table tags cover all numerical entries unless stated otherwise. Required checks are retained in the text and Appendix~\ref{app:required}; this draft is not submission-ready.}

\section{Introduction}
\label{sec:intro}
An accurate sequence classifier may rely on event identity, frequency, a recent transition, or a rule involving several positions. These explanations can produce similar test scores. Researchers need comparisons that distinguish them before attributing performance to a model's use of sequence structure.

We study this problem using sequences whose clean labels are explicitly defined. Knowing the rule makes it possible to construct inputs on which a candidate feature representation loses information and to calculate a performance reference under label noise. We compare event counts, contiguous patterns, lagged pairs and rule-specific learners. When a restricted representation predicts well, it provides an alternative explanation of the score. Determining whether a particular neural model uses that representation requires a further test of the model itself.

The tasks separate order sensitivity from information in small sets of positions. Ordered compliance (OC) is sensitive to order and has local information under its proposal distribution. Hidden-subset parity, retaining the original task abbreviation GPC, is invariant to order; its balanced version is independent of every proper fixed subset of positions. Key-inversion parity (KIP) is sensitive to order while retaining independence from sufficiently small position subsets. The constructions make these properties explicit, although differences in their input distributions prevent attributing training difficulty to a single property.

The application studies ask what remains of this approach when the label rule is unknown. In diagnosis records and listening histories, we compare predictive performance after perturbing order and examine whether the visible events are preserved. These controls evaluate the selected representation and protocol. They do not by themselves identify causal processes in disease progression or user behavior.

Our contributions are:
\begin{itemize}
\item A concrete separation of predictive performance from feature sufficiency: lag-pair prediction approaches the noisy-OC reference, while an explicit pair of inputs proves that those features do not determine the clean rule (Section~\ref{sec:synthetic}).
\item Contrasting, fully specified label rules and a parity-specific learner that establish what information the tasks contain before interpreting the outcomes of neural training (Sections~\ref{sec:rules} and~\ref{sec:synthetic}).
\item A comparison of adaptation and shuffle controls, with practical requirements for interpreting them and a clear account of the remaining evidence gaps in the model and application results (Sections~\ref{sec:setup}--\ref{sec:applications}).
\end{itemize}
The intended readers are researchers evaluating learned sequence representations. The main question is which explanation of a score survives a given comparison. Our findings support specific answers for these tasks and protocols; they do not establish a general hierarchy of model families.

\section{Related work}
\label{sec:related}
\textbf{Performance and learned rules.} Shortcut learning is a general concern in neural prediction \citep{geirhos2020}. Diagnostic examples can distinguish an intended rule from a heuristic that performs well on the usual test distribution, as illustrated by HANS \citep{mccoy2019}. Our feature comparisons follow this logic. We use the word \emph{shortcut} only for predictive features that are insufficient for the specified clean rule; simplicity or low parameter count alone does not make a solution a shortcut.

\textbf{Algorithmic sequence tasks.} Transformers can represent parity under suitable constructions \citep{chiang2022}. Expressibility, optimization, and length generalization nevertheless differ \citep{hahn2024,deletang2023}. Transformers can also learn alternative, exact algorithms for automata \citep{liu2023automata}; such algorithms need not be insufficient shortcuts in our sense. Controlled grammar studies probe learned structure in considerably richer synthetic languages \citep{allenzhu2023}. We do not claim novelty for parity or for the use of artificial sequences. Our focus is on which inference a baseline or perturbation licenses about a fixed predictor. Curriculum benefits on parity also have precedent \citep{abbe2023}, so successful length curricula must be interpreted as task-specific training results.

\textbf{Order perturbations.} Strong downstream language performance can survive shuffled pretraining \citep{sinha2021}, and recommendation datasets differ in their sensitivity to shuffling \citep{klenitskiy2024}. LLM-based recommendation has also been examined for its use of sequence information \citep{kim2025}. These findings motivate controls; they do not imply that order is generally unnecessary. Our application audit distinguishes loss of order from changes in the retained events, and recent-transition information from dependence on the entire history.

\section{Label rules and information available to a learner}
\label{sec:rules}
Let $X=(X_1,\ldots,X_n)$ be a sequence over a finite alphabet $\mathcal A$, and let $Y^*=f(X)\in\{0,1\}$ be its clean label. A hidden key set and any hidden ordering are fixed within a dataset. A \emph{task} consists of a rule, an input distribution, and a noise model.

\subsection{Ordered compliance: an ordering rule with local information}
Fix a key set $S\subseteq\mathcal A$, a rank map $\kappa:S\to\{1,\ldots,m\}$, and a lag $\lambda$. For each residue modulo $\lambda$, read the positions separated by that lag. A non-key symbol breaks the resulting chain. A \emph{maximal key run} is a consecutive stretch of keys in that chain, including every adjacent key before and after it. The ordered-compliance (OC) label is one exactly when at least one such run has length at least two and nondecreasing ranks. An ordered pair inside a longer, non-ordered run is not sufficient.

The main OC configuration is $(n,|\mathcal A|,m,\lambda)=(20,26,6,7)$, with key order \texttt{WDQJXU}. \orig\ The proposal draws symbols independently and uniformly. The original description then stratifies an oversized pool by $Y^*$ and flips labels independently with probability $\pi\in\{0,0.3\}$. \ou\ The proposal law and the distribution after selection must not be treated as identical without checking the selection procedure.

For a fixed position set $J$, write $X_J=(X_j:j\in J)$. We call the label \emph{independent of $J$} when $I(Y^*;X_J)=0$, where $I$ denotes mutual information. Independence from a few positions does not imply independence from a feature aggregate computed over the whole sequence.

\begin{proposition}[Local information in OC]
\label{prop:oc}
For the stated maximal-run rule, independent coordinates with full support, $n\geq\lambda+1$, and a nonconstant label, $I(Y^*;X_1)>0$. Independent symmetric label noise with rate below one half preserves this dependence.
\end{proposition}
The reason is concrete. When distractors exist, replace the first distractor by the smallest key while leaving the remaining positions unchanged. Existing ordered runs remain valid, possibly after extension, and some previously negative inputs become positive. Independence makes the distribution of the remaining positions the same in both cases. Appendix~\ref{app:proofs} gives the full argument, including the all-key case. The proposition establishes information at a position, not sufficient information in counts or $k$-grams, and says nothing about which features a network learns.

\subsection{Hidden-subset parity: aggregation without order}
For hidden-subset parity (GPC), let $N_a(X)$ count occurrences of symbol $a$, and fix the key set $K$. The original definition labels a sequence according to whether an even number of keys have even counts. Its equivalent, simpler form is
\begin{equation}
Y^*=c\oplus\bigoplus_{t=1}^{n}\ind\{X_t\in K\},\qquad c=(1+|K|)\bmod 2,
\label{eq:gpc}
\end{equation}
where $\oplus$ is addition modulo two. The original main setting uses $|K|=6$, alphabet size $26$, and $n=20$. \orig\ This rule is permutation-invariant: rearranging events cannot change its label.

Under independent uniform symbols, let $p=|K|/|\mathcal A|$, with $0<p<1$. The parity of the unobserved membership bits has bias $(1-2p)^{n-|J|}$. Hence $Y^*$ is independent of every proper fixed position set exactly when $p=1/2$. The balanced control uses $|K|=13$; the main six-key task is not exactly locally independent. Its one-position conditional-probability gap is $7.797207\times10^{-6}$. \newres\ The small gap does not prove that every aggregate or learning algorithm is uninformative.

\subsection{Key-inversion parity: order without local information}
The additional key-inversion parity (KIP) task fixes $S$ and $\kappa$, samples $m$ distinct positions uniformly, puts each key in exactly once in a uniformly random order, and fills the other positions independently with non-keys. Reading the keys from left to right yields ranks $(r_1,\ldots,r_m)$. Define
\begin{equation}
Y^*=\left(\sum_{i<j}\ind\{r_i>r_j\}\right)\bmod 2.
\label{eq:kip}
\end{equation}
The reported experiments use $n=20$, alphabet size $26$, $m\in\{4,6\}$ and no label noise. \ou\ This proposal is different from OC's independent-symbol proposal; its guarantees are distribution-specific.

\begin{proposition}[Order sensitivity and independence in KIP]
\label{prop:kip}
Transposing any two keys reverses the label and preserves all symbol counts. Under the KIP proposal, $Y^*$ is independent of the full count vector and of every fixed set of at most $m-2$ positions. The latter bound is tight.
\end{proposition}
Given at most $m-2$ observed positions, at least two keys remain unseen. Swapping those keys leaves the observation unchanged, preserves probability, and reverses parity. This pairs every positive completion with a negative one. Observing the first $m-1$ keys in increasing rank order can instead fix the label. The proof and a finite enumeration check appear in Appendix~\ref{app:proofs}. The construction supplies the missing order-sensitive case; it does not show that increasing local independence alone causes optimization difficulty.

\begin{table}[H]
\centering\small
\caption{The tasks separate order sensitivity from information in fixed positions. These are properties of the stated rules and proposal laws, not empirical claims about neural representations. KIP's exact bounds are established in Proposition~\ref{prop:kip}.}
\label{tab:properties}
\begin{tabular}{@{}lll@{}}\toprule
Rule & Can a permutation change $Y^*$? & Independence from fixed positions \\\midrule
OC & Yes & Fails already at $X_1$ \\
GPC, $|K|/|\mathcal A|\ne 1/2$ & No & Not exact; bias decreases with length \\
Balanced GPC & No & Every proper position subset \\
KIP & Yes & Every subset of size at most $m-2$ \\
\bottomrule\end{tabular}
\end{table}

\section{Evaluation procedure and experimental setup}
\label{sec:setup}
We first check that the stored labels match the stated rule. We then compare representations that retain different parts of the input, construct direct tests of any proposed sufficient representation, and distinguish model perturbations from retraining on perturbed data. Table~\ref{tab:comparisons} states what each comparison can establish. Source recovery and unrun model diagnostics remain required checks, rather than being treated as completed parts of the procedure.

\begin{table}[H]
\centering\small
\caption{Different comparisons answer different questions. Shuffle comparisons require the same visible events and unchanged original targets.}
\label{tab:comparisons}
\begin{tabularx}{\linewidth}{@{}>{\raggedright\arraybackslash}p{.25\linewidth}>{\raggedright\arraybackslash}X>{\raggedright\arraybackslash}X@{}}\toprule
Comparison & What it establishes & What remains open \\\midrule
Fit a restricted feature baseline & How well that representation predicts under the evaluated distribution & Whether a neural model uses those features \\
Find opposite labels with identical features & The features cannot represent the exact clean rule & The model's behavior on those inputs and their frequency \\
Evaluate a frozen model on shuffled inputs & Its sensitivity to that perturbation & Whether retraining can recover performance; effects of distribution shift \\
Train and evaluate on shuffled inputs & Performance attainable by the tested learner under the altered distribution & The algorithm learned from the original ordered inputs \\
\bottomrule\end{tabularx}
\end{table}

\textbf{Feature comparisons.} For each symbol we count occurrences. A contiguous $k$-gram scorer estimates its positive-label frequency on training data, uses $0.5$ for unseen grams, and averages scores over a test sequence. \orig\ Lag-pair features are
\begin{equation}
g_{ab}^{(d)}(X)=\sum_{t=1}^{n-d}\ind\{X_t=a,\ X_{t+d}=b\}.
\label{eq:pairs}
\end{equation}
We fit logistic regression (LR) using the supplied OC lag, and separately stack all candidate lags. Giving the lag is privileged information and is stated explicitly. These features aggregate across positions and retain some order; calling all of them ``non-sequential'' would be incorrect. On KIP, all-lag features determine pairwise key precedence when the keys are known, even though small fixed position views are independent of the label.

\textbf{Models and historical training.} The model inventory includes BERT-base \citep{devlin2019}, RoBERTa-large, Llama-3.2-1B, Llama-3.1-8B/70B, Qwen3-4B-Think, Qwen2.5-14B, a randomly initialized Llama-1M, LSTM \citep{hochreiter1997}, a small Transformer \citep{vaswani2017}, and a BiLSTM--Transformer hybrid (RNNTransformer). BiLSTM also appears in the clinical study. XGBoost \citep{chen2016} uses ordinal letter features or mean-pooled Llama embeddings. The full inventory and retained original scores are in Appendix~\ref{app:historical}.

The documented pretrained recipe uses LoRA \citep{hu2022} with rank $8$, alpha $16$, and dropout $0.1$; learning rate $2\times10^{-5}$; batch size $16$; a $200$-token cap; and at most $20$ epochs with validation-loss patience $3$. \ou\ Both BERT and the decoders use LoRA, with different target projections. Decoder classification reads the last non-padding hidden state and sums classification and language-model losses. The small Transformer uses mean pooling; recurrent models use their final state. Small-model training is documented as at most $30$ epochs with validation-F1 patience $5$. \ou\ These procedures do not isolate attention direction, readout, pretraining, or trainable parameter count. Source/manuscript disagreements are retained in Appendix~\ref{app:recipes}, rather than silently resolved into a new historical recipe.

\textbf{Data and verification.} The available OC files contain $400{,}000/50{,}000/50{,}000$ training/validation/test rows; GPC files contain $800{,}000/100{,}000/100{,}000$. \newres\ The original main table lists $400{,}000$ GPC training examples, while some training code also subsamples that dataset. \ou\ We therefore state each reanalysis's actual sample size and leave the main neural runs' exact data usage for archive verification.

The stored deterministic OC test labels disagree with the declared strict rule in $9{,}697/50{,}000$ cases ($19.394\%$). \newres\ Their positive fraction is $0.35828$, whereas the original task table reports $0.293$. \newres\ou\ This discrepancy prevents a mechanism-level interpretation of the historical deterministic results. On noisy OC, the strict rule gives latent prevalence $0.29350$ and the observed flip rate is $0.30084$ (exact $95\%$ interval $[0.29682,0.30488]$). \newres\ Those checks support the stated noisy construction. We do not relabel or overwrite historical experiments. \req{R1}{recover exact generator/run provenance or report corrected runs as new experiments; see Appendix~\ref{app:required}.}

\textbf{Metrics and uncertainty.} AUC is the probability that a positive example receives a higher score than a negative example, with half credit for ties. With latent prevalence $\rho$ and independent flip probability $\pi<1/2$, let $q=\pi+(1-2\pi)\rho$. The true-rule score has population AUC
\begin{equation}
\AUC^*=\frac12+\frac{(1-2\pi)\rho(1-\rho)}{2q(1-q)}.
\label{eq:auc}
\end{equation}
At the original stated $(\rho,\pi)=(0.293,0.3)$, it is $0.670394$. \recomp\ A fitted model can slightly exceed the empirical rule score on a finite test set without violating this population reference. F1 is threshold-dependent and the hard-rule F1 is not its maximum.

For recovered continuous predictions we compute paired DeLong AUC intervals/tests \citep{delong1988}, conditional on the fitted model and independent test rows. Six exploratory comparisons receive Holm correction \citep{holm1979}. \newres\ These quantify test uncertainty, not training-seed variation. Reported neural means/standard deviations cannot supply missing paired tests. Application predictions and most original neural runs were unavailable; uncertainty is not reconstructed from rounded tables.

\section{Synthetic results}
\label{sec:synthetic}
\subsection{Pair counts predict noisy OC but do not determine its rule}
Table~\ref{tab:ladder} shows the recomputed feature comparison. On noisy OC, the lag-pair gain over counts is $0.086996$ (paired $95\%$ interval $[0.081412,0.092580]$; Holm-adjusted $p<0.001$). \newres\ The gap between lag LR and the true-rule score is $0.000374$ ($[-0.002634,0.003381]$). \newres\ This is compatible with test variability; it is not an equivalence test or evidence that a neural model uses pairs. Removing the supplied lag is not free: the all-lag L1 fit scores lower on the noisy dataset.

\begin{table}[H]
\centering\small
\caption{Recomputed AUCs on stored data. OC test size: $50{,}000$; GPC: $100{,}000$. \newres\ Brackets are marginal $95\%$ DeLong intervals. Count/2-gram fits use the full training files; lag LR uses $200{,}000$ training rows; all-lag LR uses $100{,}000$. \newres\ $\dagger$ The stored deterministic labels disagree with the strict OC rule.}
\label{tab:ladder}
\begin{tabular}{@{}lccc@{}}\toprule
Features & Stored OC-Det.$^\dagger$ & OC-Noisy & GPC \\\midrule
Count LR \recomp & .8155 & .5839 [.5788, .5889] & .5007 \\
2-gram \recomp & .8208 & .5846 [.5795, .5896] & .4989 \\
Lag-7 LR \recomp & .9967 & .6709 [.6660, .6758] & .5024 \\
All-lag L1 LR \newres & .9973 & .6324 [.6274, .6374] & -- \\
True-rule score \newres & Not applicable & .6705 [.6665, .6745] & 1.0000 \\
\bottomrule\end{tabular}
\end{table}

The counts still fail a direct sufficiency check. Figure~\ref{fig:collision} gives opposite-label sequences with identical full symbol histograms and identical lag-pair histograms. No deterministic function of those histograms can implement the strict OC rule on every input. This statement concerns the representation, not the frequency of such examples or the computation performed by a trained model.

\begin{figure}[H]
\centering
\fbox{\begin{minipage}{.95\linewidth}\small
Let the key order be $A<B<C<D<E<F$, with $Z$ a distractor. The two length-$20$ sequences below use lag $7$. \newres
\[
\begin{array}{lll}
\texttt{ACZZZZZBBZZZZZACZZZZ} & (A,B,A),\ (C,B,C) & Y^*=0\\
\texttt{ACZZZZZBBZZZZZCAZZZZ} & (A,B,C),\ (C,B,A) & Y^*=1
\end{array}
\]
The middle column lists the maximal key runs. Both inputs have lag-pair counts $AB=BA=BC=CB=1$ and $ZZ=9$; all other pair counts are zero. \newres
\end{minipage}}
\caption{An exact diagnostic example. One input contains an ordered maximal run and the other does not, despite identical count and lag-pair features. Symbol names are illustrative rank-preserving relabelings.}
\label{fig:collision}
\end{figure}

\req{R2}{evaluate frozen neural models on diverse pairs of this type, with opposite clean labels and identical feature counts. Report within-pair ranking and uncertainty over seeds and template families; estimated 1--4 H200-hours for evaluation.} Failure would refute exact rule recovery on these inputs; successful discrimination would refute a pair-count-only account for those models. Neither result alone identifies their internal algorithm.

On the noisy test set, the hard-rule F1 is $0.578396$, while predicting positive for every example yields $0.591609$. \newres\ Similar F1 values are therefore not evidence that two models learned the same ordering information. We retain AUC and validation-selected threshold metrics separately.

\subsection{Parity failure is specific to a learning procedure}
The reported main neural runs have near-chance GPC AUCs across the tested neural families, including a single $70$B decoder run. \ou\ The reported binary-parity and curriculum results rule out a universal claim that these architectures cannot learn parity. In the reported length curriculum, LSTM succeeds in all three seeds and Transformer in two of three. \ou\ The exact $95\%$ success-probability intervals from those counts are $[0.292,1]$ and $[0.094,0.992]$, respectively. \newres\ Three runs do not establish reliable success rates or isolate why a run failed.

The hidden subset is learnable with a suitable hypothesis class. Set $z_a=N_a(X)\bmod2$ and let $w_a$ indicate membership in $K$. Equation~\ref{eq:gpc} gives linear equations in $w$ over the binary field. Gaussian elimination on the first $256$ training examples yields rank $25$ and complementary solutions of sizes six and twenty; the supplied subset size selects the six-key solution. \newres\ Validation and test each have zero errors on $100{,}000$ examples and AUC $1$. \newres\ The exact two-sided $95\%$ upper bound on test error is $0.00003689$. \newres

This learner knows the parity family but is not given the key identities. It shows that data-level subset identification is possible, not that a neural optimizer will discover the same computation. The reported membership-revealed multilayer perceptron (MLP) controls likewise change the information supplied to the learner; they cannot uniquely locate a neural bottleneck. \req{R3}{verify the original curriculum archive and compare longer direct/curriculum training at matched optimizer steps, without early stopping; the specified five-seed small-model study is estimated at 30--50 H200-hours.}

\subsection{Full adaptation changes the reported decoder comparison}
\label{sec:adaptation}
The additional results report full-backbone fine-tuning on noisy OC, with the input, classification readout, loss, learning rate and schedule otherwise held fixed. Table~\ref{tab:adaptation} retains these values alongside the original LoRA results. The reported improvement contradicts an interpretation of the earlier plateau as an intrinsic inability of causal decoders. It does not show that full fine-tuning is necessary, that a different LoRA configuration could not succeed, or that the successful models recovered the strict rule.

\begin{table}[H]
\centering\small
\caption{Reported noisy-OC AUCs; all entries are \ou. Parentheses in the original rows are reported standard deviations. No seed count or interval accompanies the new full-fine-tuning values. Total model size is not the number of updated parameters.}
\label{tab:adaptation}
\begin{tabular}{@{}lll@{}}\toprule
Model & Adaptation & AUC \\\midrule
BERT-base & Rank-8 LoRA & .669 (.006) \\
Llama-3.2-1B & Rank-8 LoRA & .587 (.008) \\
Llama-3.2-1B & Full fine-tuning & .6703 \\
Llama-3.1-8B & Rank-8 LoRA & .575 (.008) \\
Llama-3.1-8B & Full fine-tuning & .6684 \\
Qwen2.5-14B & Rank-8 LoRA; single run & .670 \\
\bottomrule\end{tabular}
\end{table}

\req{R4}{recover configurations, exact trainable counts, checkpoints and continuous predictions for the full-fine-tuning runs; verify the stated matching and add seed/test uncertainty before claiming adaptation efficiency. Archive-based re-evaluation is estimated at 0--4 H200-hours; a training rerun needs a measured pilot.} The original objective, readout and tuning differences still prevent an architecture-only comparison. The reported Qwen success changes both scale and family, so it is not a controlled scaling result.

\subsection{Inversion parity adds a positive order-learning case}
Table~\ref{tab:kip} summarizes the additional KIP results. The construction itself proves count independence. The model controls address the separate question in Table~\ref{tab:comparisons}: whether successful prediction depends on order. Training on shuffled inputs with unchanged original labels tests for a surviving order-invariant route; evaluating a frozen ordered model on shuffled inputs tests its sensitivity to that perturbation. These controls are not interchangeable.

\begin{table}[H]
\centering\small
\caption{Reported LSTM results on KIP with four keys; every numerical entry is \ou. Three ordered training seeds are reported. The supplied $\pm$ quantities have not been identified as standard deviations or confidence intervals.}
\label{tab:kip}
\begin{tabular}{@{}lll@{}}\toprule
Training input & Test input & Reported AUC \\\midrule
Ordered & Ordered & 1.000 in each seed \\
Shuffled, original labels & Shuffled, original labels & $0.500\pm0.002$ \\
Ordered, frozen checkpoint & Shuffled, original labels & $0.502\pm0.006$ \\
\bottomrule\end{tabular}
\end{table}

The report describes near-chance four-key results for the small Transformer, BERT-LoRA, Llama-1B-LoRA and full-fine-tuned decoder controls; exact per-run values are missing. With six keys, it reports failure of all tested raw-sequence models, including LSTM. \ou\ Increasing the key count changes the inversion computation as well as the independence threshold, so this comparison cannot attribute difficulty to local independence alone. \req{R5}{verify the KIP generator, split hashes, complete run matrix, shuffle transformations and per-seed predictions; estimated 0--4 H200-hours for checkpoint evaluation, with any retraining budget set by a pilot.} The mathematical guarantee is established here; the positive neural demonstration remains conditional on that verification.

\section{Applications: distinguish order from retained content}
\label{sec:applications}
\subsection{Clinical diagnosis sequences}
We retain the original MIMIC-IV cohort \citep{johnson2023} for chronic kidney disease (CKD) and end-stage renal disease (ESRD): $14{,}813$ patients, $1{,}146$ positives and $293$ Clinical Classifications Software (CCS) diagnosis categories. \ou\ Positives contain records before the first ESRD-related admission; negatives contain their available histories. Same-admission CKD/ESRD cases are excluded. Codes are ordered by admission time and, within admissions, by diagnosis sequence number. That within-admission rank is not an event timestamp. Appendix~\ref{app:clinical} gives the unchanged endpoint codes and exclusions.

This is a retrospective diagnosis-sequence audit. The cohort has no common prediction landmark and fixed future horizon, so it does not establish prospective CKD progression performance. Table~\ref{tab:clinical} shows points verified against aggregate training logs. The original rounded values are retained in the supplement. Bag-of-codes LR and the reported unigram scorer also perform well, but use different summary/split evidence and do not resolve the neural order comparison.

\begin{table}[H]
\centering\small
\caption{Clinical AUCs verified from aggregate logs \orig. Each condition trains and evaluates a separate model on the same patient split. Differences are arithmetic from those points \newres; no paired patient or training-seed intervals are available. Shuffling before truncation confounds order with retained content.}
\label{tab:clinical}
\begin{tabular}{@{}lrrr@{}}\toprule
Model & Ordered & Shuffled & Difference \\\midrule
Transformer & .8596 & .8436 & +.0160 \\
BiLSTM & .8203 & .8227 & -.0024 \\
BERT & .8258 & .8426 & -.0168 \\
LSTM & .5000 & .4672 & +.0328 \\
\bottomrule\end{tabular}
\end{table}

Two implementation issues limit interpretation. The pipeline shuffles full histories before model-specific truncation, so long ordered/shuffled inputs can contain different codes. The recurrent implementation reads states after padding, which provides an alternative explanation for its failed LSTM result. The observed differences therefore do not establish that diagnosis order is unnecessary.

\req{R6}{retain the same codes before shuffling; correct recurrent length handling; repeat Transformer/LSTM/BiLSTM over five training seeds and three shuffle draws; report average precision and paired patient/seed uncertainty. Estimated 4--12 H200-hours, inside the credentialed environment.} A bounded ``small order effect'' claim requires a prespecified practical margin and a sufficiently narrow interval, not merely a nonsignificant test.

\subsection{Next-song prediction}
The additional 30Music experiment \citep{turrin2015} compares SASRec \citep{kang2018}, GRU4Rec \citep{hidasi2015}, and simple recommendation baselines. HR@10 is the fraction of targets appearing among the ten highest-ranked candidates. The reported main protocol retains the same final $128$ songs before shuffling. \ou\ ``Last fixed'' preserves the last song; ``recent half fixed'' preserves the later half of the retained history. Table~\ref{tab:music} reports all supplied values.

\begin{table}[H]
\centering\small
\caption{Reported 30Music HR@10 under a fixed visible-song multiset. All entries: \ou. The report mentions three training seeds and three shuffle draws, with variation at most $0.005$; this is not a defined uncertainty interval.}
\label{tab:music}
\begin{tabular}{@{}lrrrr@{}}\toprule
Predictor & Ordered & All shuffled & Last fixed & Recent half fixed \\\midrule
SASRec & .303 & .041 & .255 & .303 \\
GRU4Rec & .278 & .029 & .099 & .274 \\
Trigram with bigram fallback & .309 & .019 & .300 & .309 \\
Bigram & .300 & .019 & .300 & .300 \\
Item similarity & .065 & .065 & .065 & .065 \\
Most popular & .003 & .003 & .003 & .003 \\
\bottomrule\end{tabular}
\end{table}

The reported pattern differs from the clinical results: removing order substantially lowers prediction, while recent transitions account for much of the useful information. The bigram/trigram scores match or exceed the neural points; this does not demonstrate dependence on a long history. The report also reproduces a separate shuffle-before-truncation protocol at $0.197\to0.020$, against published $0.198\to0.020$. \ou\ Those numbers must not be compared directly with the fixed-content table because the protocols differ.

\req{R7}{recover preprocessing, temporal splits, candidate/negative sampling, training-versus-test shuffle behavior and per-user predictions; verify train/test separation and equal retained songs. Report paired intervals across users, training seeds and shuffle draws. Estimated 0--4 H200-hours for archived-model scoring; retraining requires a pilot.} Until then, the table is an additional reported application, not an independently reproduced validation.

\section{Discussion}
The OC analysis shows why a strong feature baseline should lead to a model diagnostic. Near-reference AUC is attainable with pair counts, but the same pair counts can correspond to opposite clean labels. Testing neural predictions on those inputs would distinguish a pair-count-only explanation from alternatives. Agreement on ordinary test AUC cannot make that distinction. The diagnostic remains unfinished for the trained networks, and the current evidence establishes insufficiency of the representation rather than the internal computation of a model.

The parity comparisons make a related distinction between information and optimization. A known hypothesis class can make the hidden subset easy to identify, whereas a neural training procedure may not find that computation. Conversely, KIP requires order without giving label information to a few fixed positions. Its all-lag representation can still retain the needed information. These examples show why independence from a restricted observation should not be promoted into a learnability claim about full sequences.

The reported fine-tuning results change the interpretation of the original decoder gap: adaptation choices remain part of the explanation. In applications, a shuffle-related decline also admits more than one interpretation. It demonstrates sensitivity to the perturbation, potentially including distribution shift. Retrained controls and preservation of recent context address different alternatives. Retaining the same events is necessary to separate these effects from a change in available content.

The practical procedure is therefore to state the candidate explanation, specify a comparison that could contradict it, and report the limits of that comparison. A successful predictor, a sufficient representation and a learned algorithm each require their own evidence. This principle is useful beyond the present tasks, while the empirical findings remain limited to the stated datasets, training procedures and available outputs.

\section{Limitations}
\label{sec:limitations}
The synthetic rules are deliberately simple and do not approximate natural-language meaning or clinical causation. GPC is an order-invariant control; OC and KIP use different proposal laws. Generalization across hidden sets, ranks and lengths is limited. The deterministic OC provenance discrepancy is unresolved. Most original neural metrics and the new full-fine-tuning, KIP and 30Music findings lack available run archives; their tags cannot be removed without verification. Existing baselines establish representational sufficiency or insufficiency, not internal neural mechanisms. Single large-model runs and small seed counts do not support scale laws or precise success probabilities. The clinical construction and truncation issue prevent prospective or general clinical-order claims. An additional patient-level analysis must remain within credentialed access. These limitations constrain the conclusions even if all reported point estimates are reproduced.

\section{Conclusion}
Known label rules make sequence prediction auditable: they let us distinguish accurate scores from sufficient features and successful training procedures. OC, hidden-subset parity and inversion parity separate order sensitivity from local information. The resulting controls support specific statements about model behavior while keeping adaptation, data integrity and application protocol part of the explanation.
\FloatBarrier
\label{LastMainPage}
\subsubsection*{Reproducibility statement}
Definitions and proofs specify the synthetic rules. The feature reanalysis records input hashes, estimator settings, predictions and statistical computations. Appendices provide historical recipes, original results and explicit verification requirements; unverified values are not represented as reproduced. The anonymous repository link is retained: \url{https://anonymous.4open.science/r/anonymous_neurips_seqllm/}. No new training is claimed in this revision. The repository still needs the missing run archives, corrected experiment versions and final figure-generation material before submission.

\subsubsection*{Broader impact and ethics statement}
MIMIC-IV requires credentialed access, training and a data-use agreement \citep{johnson2023}. The clinical audit uses the original cohort definition and is intended to examine representation-dependent prediction, not to provide a clinical model or patient-care recommendation. Patient-level data, derived code sequences, identifiers and predictions are not released. Only approved aggregate outputs and analysis code may leave the credentialed environment. Misinterpreting predictive scores as evidence of clinically meaningful temporal reasoning could encourage unjustified confidence in a model; the stated limits and fixed-content checks address that risk. The revision's clinical checks used cohort scripts and aggregate logs, without accessing patient records. The synthetic labels are generated by explicit rules; language models do not annotate them.

\subsubsection*{AI assistance disclosure}
The original work disclosed AI assistance with code debugging, grammar improvement and visualization scripts. During this revision, an AI assistant additionally helped review the manuscript and reviewer responses, inspect code and aggregate results, perform and document reanalyses, check mathematical examples and references, interpret findings, plan controls, draft and restructure the text, and prepare the journal-format revision. These activities extend beyond language editing. No AI-generated empirical result substitutes for an executed experiment; unrecovered results and unrun controls are explicitly marked. The authors remain responsible for reviewing the final text, proofs, references and numerical provenance. \req{R8}{authors confirm and complete this disclosure, including any earlier assistance not recorded here.}

\bibliography{references}
\bibliographystyle{tmlr}
\FloatBarrier
\appendix
\section{Proofs and direct checks}
\label{app:proofs}
\subsection{Local information in the maximal-run OC rule}
For each fixed value $a$ at the first position, let $F_a(x_2,\ldots,x_n)$ be the binary OC label. Independence means that the same distribution of $(X_2,\ldots,X_n)$ is used when comparing $F_a$ and $F_b$. This is all that is meant by coupling the two inputs: we draw the remaining positions once and compare two choices at the first position. Identical coordinate distributions are not required.

Suppose first that $S\ne\mathcal A$. Let $s_{\min}$ be the minimum-rank key and $a_0$ a distractor. Any ordered maximal run under $X_1=a_0$ either remains unchanged under $X_1=s_{\min}$ or starts at $1+\lambda$ and extends to position one. Such an extension is still nondecreasing because its new first rank is the minimum. Thus $F_{s_{\min}}\geq F_{a_0}$ pointwise. On the event that $X_{1+\lambda}\in S$ and every other position after the first is a distractor, the first choice gives an ordered two-key run and the second gives no ordered run. Full support makes this event have positive probability. Consequently the two conditional label probabilities differ.

When $S=\mathcal A$, every residue class modulo $\lambda$ is a maximal run. A run of length one cannot certify the label. Compare the minimum and maximum key at position one. Nondecreasing runs possible under the maximum are also possible under the minimum. The inclusion is strict: assigning the minimum to the rest of the first residue run produces an ordered run only under the minimum, provided distinct ranks exist. Independently, every other residue run of length at least two can be made non-ordered by starting with a maximum and following it with a minimum. Full support gives this joint event positive probability. Other residue runs therefore cannot erase the strict conditional difference. The nonconstant-label assumption excludes the degenerate single-rank case.

For independent label flips, the observed conditional probability is
\[
P(Y=1\mid X_1=a)=\pi+(1-2\pi)P(Y^*=1\mid X_1=a).
\]
For $\pi<1/2$ this map is strictly increasing, so the observed label also depends on the first position. Under uniform symbols with a proper key set, the witness event yields the original bound
\[
P(Y^*=1\mid X_1=s_{\min})-P(Y^*=1\mid X_1=a_0)
\geq \frac{m}{|\mathcal A|}\left(\frac{|\mathcal A|-m}{|\mathcal A|}\right)^{n-2}.
\]
Noise multiplies that difference by $1-2\pi$. This is a lower bound on a conditional probability difference; it is not a guarantee that a selected feature family attains the Bayes score. The result concerns this OC rule and independent full-support inputs, not all order-based tasks or dependent clinical records.

\subsection{Parity and the noise reference}
The parity of the number of even-count keys equals $(|K|+\sum_{a\in K}N_a)\bmod2$. Converting ``even'' into label one gives the fixed flip in Equation~\ref{eq:gpc}. Under independent uniform symbols, $B_t=\ind\{X_t\in K\}$ are independent Bernoulli($p$), so
\[
\mathbb E\!\left[(-1)^{\sum_{t\notin J}B_t}\right]=(1-2p)^{n-|J|}.
\]
At $p=1/2$, every nonempty unobserved set has unbiased parity, independently of $X_J$. Away from that balance, the bias is nonzero and the observed parity can change the conditional label distribution. For fixed $J$ and $0<p<1$, its magnitude tends to zero with increasing length. Rearranging positions never changes the full parity. The theorem is therefore about information in a restricted observation, not order sensitivity or impossibility for a learner that sees all positions.

For Equation~\ref{eq:auc}, put $a=P(Y^*=1\mid Y=1)=\rho(1-\pi)/q$ and $b=P(Y^*=1\mid Y=0)=\rho\pi/(1-q)$. The binary true-rule score has AUC $(1+a-b)/2$, giving the stated formula. Its hard-classifier precision, recall and F1 are
\[
1-\pi,\qquad \frac{\rho(1-\pi)}{q},\qquad \frac{2\rho(1-\pi)}{\rho+q}.
\]
These are properties of that classifier. Optimizing F1 over thresholds can prefer the all-positive classifier, so the hard-rule F1 is not an upper bound.

\subsection{Inversion parity and feature sufficiency}
Condition on any attainable observation $X_J$ with $|J|\leq m-2$. At least two keys have not been observed. Choose two of those key identities by a fixed rule and exchange their positions in every completion. The operation is its own inverse, preserves the key-position set and all distractors, and is probability-preserving under the uniform key assignment. A transposition reverses permutation parity, so the two labels have equal conditional probability. The same argument, conditional on the full count vector, proves count independence.

For tightness, take $J=\{1,\ldots,m-1\}$ and condition on those positions containing ranks $1,\ldots,m-1$ in order. This event has positive probability. The final key must occur later, so the key order is increasing and the label is fixed. Independence thus fails for at least one view of size $m-1$. With four keys, independence holds through two positions; with six, through four. \newres

An exhaustive illustrative check with four keys, length six and one distractor enumerated $360$ sequences and all $22$ position sets of sizes zero, one or two. Every conditional label distribution was balanced, and the three-position witness fixed the label. \newres\ This checks the finite example, not the unverified training experiments or every possible parameter setting.

There are two useful cautions. First, because each key occurs exactly once, changing its fixed rank map changes the inversion label only by a dataset-wide bit flip: permutation signs multiply. A classifier need not identify the entire hidden ranking. Second, $\sum_d g_{ab}^{(d)}$ records whether key $a$ precedes key $b$, so all-lag counts determine inversion parity when key identities and ranking are supplied. This is a global representation, not a contradiction of independence from a small fixed set of positions. A learned readout on that representation can still fail to find parity.

\section{Recipes, provenance and reanalysis details}
\label{app:recipes}
\textbf{Historical pretrained recipe.} The documented Llama/Qwen classifier receives \texttt{Sequential events:} followed by space-separated letters and the sentinel \texttt{Outcome (0 or 1):}; the true label is not appended. The classifier is a $d\to d/2\to2$ head with Tanh and dropout $0.1$, reading the last non-padding hidden state. Its loss is classification cross-entropy plus autoregressive prompt loss with equal weights. LoRA targets q/k/v/o projections; BERT's targets query/key/value projections and uses its sequence-classification head. These are reported settings, not newly inferred historical configurations. \ou

The manuscript specifies AdamW, $6\%$ warmup followed by linear decay, learning rate $2\times10^{-5}$, batch size $16$, token cap $200$, LoRA rank/alpha $8/16$, adapter dropout $0.1$, and loss-based early stopping with patience $3$ and maximum $20$ epochs. \ou\ Some saved sensitivity jobs instead specify BERT batch size $32$ and Llama batch size $64$ with token cap $100$. \orig\ Those logs are not proof of the main table's exact configuration. Total and trainable parameter counts must be recovered separately for every compared adaptation condition (R4).

\begin{table}[H]
\centering\small
\caption{Reported small-model settings. All numerical entries: \ou. FFN denotes the Transformer feed-forward width. The two configurations are retained as reported; the missing run manifests prevent mapping every score to its exact configuration.}
\begin{tabular}{@{}llrrrr@{}}\toprule
Model & Configuration & Embedding & Hidden/FFN & Layers & Learning rate \\\midrule
Transformer & Default & 64 & 256 & 2 & $5\times10^{-4}$ \\
Transformer & Optimized & 64 & 128 & 3 & $5\times10^{-4}$ \\
LSTM & Default & 32 & 64 & 2 & $10^{-3}$ \\
LSTM & Optimized & 32 & 128 & 2 & $2\times10^{-3}$ \\
RNNTransformer & Default & 64 & 64/256 & 1+2 & $5\times10^{-4}$ \\
RNNTransformer & Optimized & 128 & 32/256 & 1+3 & $10^{-4}$ \\
\bottomrule\end{tabular}
\end{table}

The small Transformer uses four attention heads, GELU, learned position embeddings and mean pooling with a padding mask. \ou\ Default small-model dropout is $0.3$; the optimized LSTM uses $0.2$. Batch size is $64$ and the reported schedule permits $30$ epochs, with validation-F1 patience $5$. The randomly initialized Llama-1M has hidden width $128$, four layers, four attention heads, two key/value heads and intermediate width $512$. \ou\ The manuscript calls the optimizer AdamW, while the inspected small-model trainer instantiates Adam. It describes letters indexed from one, while the inspected dataset maps A to zero and the Transformer treats zero as padding. The trainer's shallow best-state dictionary copy also retains references to changing tensors. \newres\ These findings require source-version recovery or clearly labeled corrected runs; they do not quantify the change in any historical score.

\textbf{XGBoost and the feature refits.} The documented main XGBoost settings are depth $6$, minimum child weight $2$, split penalty $0.1$, row/column subsampling $0.8$, L1/L2 regularization $0.5/1.5$, learning rate $0.05$, at most $200$ rounds and validation-log-loss patience $10$. \ou\ The embedding variant uses mean-pooled Llama-1B features. The feature-audit table's earlier ``BoE-XGB'' label instead refers to scikit-learn gradient boosting, with $200$ estimators, depth $4$, subsample $0.8$ and seed $42$. Its reproduced AUCs are $0.815804$, $0.584070$ and $0.502810$ on stored deterministic OC, noisy OC and GPC. \recomp\ This naming correction does not replace the distinct main XGBoost models.

Count and supplied-lag LR use L2 regularization with $C=1$; count features cover the full alphabet. The all-lag L1 fit uses $C=0.5$. \newres\ The original key-count feature list used WDQJXN, whereas the verified noisy OC key order ends in U. Replacing N by U gives noisy key-count AUC $0.583681$, compared with $0.574264$ as coded. \newres\ This correction is separated from the full-alphabet count baseline. The original deterministic generator remains unidentified.

\textbf{Splits, selection and statistical scope.} The original neural experiment description reports training fractions $1\%,10\%,30\%,50\%,75\%,100\%$, generally three seeds, and threshold selection over $0.10,0.11,\ldots,0.89$ to maximize validation F1. It also names seed $9550$ without resolving the full run list. \ou\ Qwen-14B and Llama-70B rows explicitly report single runs. The claimed original hardware is one H100 with $80$GB memory. \ou\ The H200 estimates in the TODOs concern future checks, not the hardware of those historical results.

The reanalysis uses the supplied split files, scikit-learn $1.5.1$, seed $42$ where applicable, DeLong's classwise-placement variance with half credit for ties, and paired differences on identical test rows. \newres\ The six exploratory comparisons are lag versus counts in all three datasets, corrected versus coded key counts in the two OC datasets, and lag versus the true-rule score on noisy OC. \newres\ Their adjusted p-values do not identify neural mechanisms. Main neural learning curves, most per-seed predictions and the exact main run configurations were not recovered. They are not replaced by curves drawn from hardcoded table values. Saved count/$k$-gram/boosting outputs reproduce in all $18$ entries to maximum absolute AUC discrepancy $1.11\times10^{-16}$. \recomp

\section{Historical results retained without numerical changes}
\label{app:historical}
The following compact tables preserve the original main performance table, including metrics not emphasized in the revised argument. Asterisks identify the two stated single runs; the dagger identifies the BiLSTM--Transformer hybrid. A printed zero standard deviation or perfect mean is retained as supplied, not promoted to certainty. Reported model sizes are BERT $110$M, RoBERTa $355$M, the named decoder sizes, LSTM $198$K--$340$K, BiLSTM $256$K--$395$K, Transformer $103$K--$104$K, and the hybrid $319$K--$455$K. \ou\ Updated-parameter counts are unavailable. The below-chance noisy-OC XGBoost points require score-orientation and split checks before interpretation; their values are retained, without claiming systematic anti-prediction.

The Naive rule checks whether the full sequence is alphabetically nondecreasing. In the available test file, $1{,}953/40{,}000$ sorted sequences are labeled negative, and the strict-rule score has AUC $0.951224$, rather than one. \newres\ Its historical table therefore serves as an unresolved sanity check. The deterministic OC caveat in Section~\ref{sec:setup} applies to every corresponding historical row. The historical GPC 1-gram AUC $.505$ does not match the supplied scorer's recomputed $.500652$. \ou\recomp\ These discrepancies are preserved explicitly.

\begin{table}[H]
\centering\small
\caption{Historical results: Naive ($\pi=0.0$). Values are transcribed without numerical changes from the supplied manuscript; parentheses are its reported standard deviations. All numerical entries: \ou. Run-level verification remains incomplete.}
\begin{tabular}{@{}lcccc@{}}\toprule
Model & AUC & F1 & Precision & Recall \\\midrule
\textit{Oracle} & 1.000 & 1.000 & 1.000 & 1.000 \\
$k$-gram & .999 & .999 & .999 & .999 \\
XGBoost (basic) & .603 (.004) & .714 (.005) & .556 (.004) & .995 (.003) \\
XGBoost (LLM) & .502 (.004) & .536 (.008) & .501 (.004) & .577 (.010) \\
BERT & .999 (.001) & .999 (.001) & .999 (.001) & .999 (.001) \\
RoBERTa & .999 (.001) & .999 (.001) & .999 (.001) & .999 (.001) \\
Llama-3.2-1B & .999 (.001) & .999 (.001) & .999 (.001) & .999 (.001) \\
Llama-3.1-8B & .999 (.001) & .999 (.001) & .999 (.001) & .999 (.001) \\
Llama-1M (cold) & .999 (.001) & .999 (.001) & .999 (.001) & .999 (.001) \\
Qwen3-4B-Think & .999 (.001) & .999 (.001) & .999 (.001) & .999 (.001) \\
LSTM & .999 (.001) & .999 (.001) & .999 (.001) & .999 (.001) \\
Transformer & .999 (.001) & .999 (.001) & .999 (.001) & .999 (.001) \\
RNNTransformer$^{\dagger}$ & .999 (.001) & .999 (.001) & .999 (.001) & .999 (.001) \\
\bottomrule\end{tabular}
\end{table}
\begin{table}[H]
\centering\small
\caption{Historical results: OC-Deterministic ($\pi=0.0$). Values are transcribed without numerical changes from the supplied manuscript; parentheses are its reported standard deviations. All numerical entries: \ou. Run-level verification remains incomplete.}
\begin{tabular}{@{}lcccc@{}}\toprule
Model & AUC & F1 & Precision & Recall \\\midrule
\textit{Oracle} & 1.000 & 1.000 & 1.000 & 1.000 \\
2-gram & .821 & .672 & .593 & .775 \\
XGBoost (basic) & .625 (.012) & .457 (.015) & .516 (.010) & .409 (.018) \\
XGBoost (LLM) & .494 (.010) & .399 (.012) & .353 (.008) & .459 (.015) \\
BERT & .999 (.001) & .999 (.001) & .999 (.001) & 1.000 (.001) \\
RoBERTa & .999 (.001) & .999 (.001) & .999 (.001) & 1.000 (.001) \\
Llama-3.2-1B & .846 (.015) & .695 (.018) & .623 (.020) & .785 (.022) \\
Llama-3.1-8B & .995 (.003) & .984 (.008) & .976 (.010) & .993 (.007) \\
Llama-1M (cold) & .832 (.012) & .686 (.015) & .605 (.018) & .791 (.020) \\
Qwen3-4B-Think & .850 (.014) & .701 (.016) & .620 (.017) & .800 (.020) \\
LSTM & .999 (.001) & .994 (.003) & .995 (.003) & .994 (.003) \\
Transformer & .999 (.001) & .992 (.002) & .985 (.005) & 1.000 (.001) \\
RNNTransformer$^{\dagger}$ & .997 (.002) & .992 (.003) & .985 (.005) & 1.000 (.002) \\
\bottomrule\end{tabular}
\end{table}
\begin{table}[H]
\centering\small
\caption{Historical results: OC-Noisy ($\pi=0.3$). Values are transcribed without numerical changes from the supplied manuscript; parentheses are its reported standard deviations. All numerical entries: \ou. Run-level verification remains incomplete.}
\begin{tabular}{@{}lcccc@{}}\toprule
Model & AUC & F1 & Precision & Recall \\\midrule
\textit{Oracle} & .670 & .577 & .700 & .491 \\
2-gram & .585 & .592 & .420 & 1.000 \\
XGBoost (basic) & .463 (.010) & .418 (.012) & .437 (.008) & .400 (.015) \\
XGBoost (LLM) & .418 (.011) & .453 (.014) & .437 (.008) & .470 (.018) \\
BERT & .669 (.006) & .592 (.004) & .420 (.005) & 1.000 (.002) \\
RoBERTa & .652 (.008) & .592 (.003) & .420 (.005) & 1.000 (.002) \\
Llama-3.2-1B & .587 (.008) & .592 (.005) & .420 (.006) & .999 (.005) \\
Llama-3.1-8B & .575 (.008) & .591 (.005) & .425 (.006) & .968 (.010) \\
Qwen3-4B-Think & .597 (.007) & .591 (.004) & .423 (.005) & .979 (.008) \\
Qwen2.5-14B$^{*}$ & .670 (---) & .586 (---) & .476 (---) & .763 (---) \\
LSTM & .665 (.003) & .562 (.020) & .660 (.002) & .489 (.031) \\
Transformer & .671 (.001) & .579 (.003) & .701 (.015) & .492 (.012) \\
RNNTransformer$^{\dagger}$ & .671 (.005) & .578 (.005) & .683 (.010) & .501 (.005) \\
\bottomrule\end{tabular}
\end{table}
\begin{table}[H]
\centering\small
\caption{Historical results: GPC-Deterministic ($\pi=0.0$). Values are transcribed without numerical changes from the supplied manuscript; parentheses are its reported standard deviations. All numerical entries: \ou. Run-level verification remains incomplete.}
\begin{tabular}{@{}lcccc@{}}\toprule
Model & AUC & F1 & Precision & Recall \\\midrule
\textit{Oracle} & 1.000 & 1.000 & 1.000 & 1.000 \\
1-gram & .505 & .668 & .501 & 1.000 \\
XGBoost (basic) & .495 (.004) & .490 (.005) & .497 (.003) & .484 (.006) \\
XGBoost (LLM) & .501 (.003) & .639 (.005) & .502 (.002) & .881 (.008) \\
BERT & .503 (.002) & .668 (.002) & .501 (.002) & 1.000 (.001) \\
RoBERTa & .501 (.003) & .668 (.002) & .501 (.002) & 1.000 (.001) \\
Llama-3.2-1B & .497 (.005) & .668 (.001) & .501 (.002) & 1.000 (.001) \\
Llama-3.1-8B & .500 (.003) & .668 (.002) & .501 (.002) & 1.000 (.001) \\
Llama-1M (cold) & .503 (.003) & .668 (.002) & .501 (.003) & 1.000 (.001) \\
Llama-3.1-70B$^{*}$ & .502 (---) & .668 (---) & .501 (---) & 1.000 (---) \\
Qwen3-4B-Think & .500 (.002) & .668 (.001) & .501 (.000) & 1.000 (.001) \\
LSTM & .502 (.003) & .668 (.002) & .501 (.002) & 1.000 (.001) \\
Transformer & .501 (.003) & .668 (.002) & .501 (.002) & 1.000 (.001) \\
RNNTransformer$^{\dagger}$ & .501 (.002) & .668 (.002) & .501 (.003) & 1.000 (.001) \\
\bottomrule\end{tabular}
\end{table}

\FloatBarrier

\textbf{Additional historical controls.} The reported binary-parity reference has $7/10$ successful initializations at length ten and none at length twenty; the same failing initializations recur across fixed and streaming data, so those regimes cannot be pooled as independent seeds. \ou\ The original length-curriculum results at the final length, in seed order $0,1,2$, are Transformer $(0.999,0.497,0.999)$, LSTM $(1.000,1.000,1.000)$ and BERT $(0.493,0.502,0.500)$. \ou\ These support the stated success counts, with considerable uncertainty. The sensitivity sweep's recovered logs contain one seed per configuration, not a multi-seed uncertainty estimate. The held-out-rule script uses noiseless ordered-pair existence rather than strict noisy OC and selects an F1 threshold on held-out labels; it does not verify transfer of the rule defined here. Attention plots and few-shot summaries are excluded from the main evidence because their exact run/input provenance is missing. Original table blocks remain unchanged in the editorial evidence archive rather than being reproduced as lengthy, load-bearing appendices.

\section{Application details and additional reported controls}
\label{app:clinical}
\textbf{Unchanged clinical cohort.} The source describes MIMIC-IV v3.1, with CKD codes ICD-10 N18.1--N18.5 and ICD-9 585.1--585.5. ESRD-related codes are ICD-10 N18.6, Z99.2, Z49.* and ICD-9 585.6, V45.1, V56.*. \orig\ Positives have a first ESRD-related admission strictly after the first CKD admission, with sequences ending before that ESRD admission. Negatives have CKD without a recorded ESRD-related endpoint. Same-admission cases and records with fewer than ten codes after processing are excluded. \orig\ Because histories are selected differently by outcome and no shared landmark/horizon is fixed, the label should not be described as prospective progression at a common forecast time.

Reported sequence lengths are median $49$, mean $82$ and maximum $2{,}396$, with positive/negative medians $54/49$. The cohort contains $13{,}667$ negatives and positive prevalence $7.7\%$. \ou\ The original neural/bag-of-codes split is $60/20/20$, whereas $k$-gram summaries average five $80/20$ splits. \ou\ The code prepares one shuffled history per patient before truncation; small models cap at $1{,}024$ codes, and BERT truncates to $512$ tokens. \orig\ These differing retained inputs must be checked inside the credentialed environment. The recurrent models do not pack sequences in the inspected implementation, so reading their final hidden state includes padding. The clinical $k$-gram script selects its F1 threshold on test labels, invalidating that threshold-based evaluation, but not its AUC solely for this reason.

\begin{table}[H]
\centering\small
\caption{Original clinical summary values retained verbatim. All entries: \ou. These summaries are not matched paired estimates: splits differ as stated above, and the $k$-gram F1 thresholds were selected on test data.}
\begin{tabular}{@{}lrrl@{}}\toprule
Model/control & Ordered AUC & Shuffled AUC & Original auxiliary value \\\midrule
1-gram & $.823\pm.015$ & -- & F1 .356 \\
2-gram & $.801\pm.016$ & -- & F1 .341 \\
3-gram & $.576\pm.015$ & -- & F1 .165 \\
4--7 gram & .479--.500 & -- & F1 .144--.146 \\
Bag-of-codes LR & .809 & .809 & Difference 0 \\
Bag-of-codes boosting & .803 & .803 & Difference 0 \\
Transformer & .860 & .844 & Difference +.016 \\
BiLSTM & .820 & .823 & Difference -.003 \\
BERT & .826 & .843 & Difference -.017 \\
LSTM & .500 & .467 & Failed run \\
\bottomrule\end{tabular}
\end{table}

\textbf{Additional KIP details.} The supplied report gives bag-of-letter AUCs $0.498$--$0.509$, $k$-gram AUCs $0.499$--$0.503$, and fixed-lag baseline AUCs at most $0.524$ for the four-key setting. It reports all-lag performance of $1.0$ there, but does not recover the exact fitted readout in the available materials. For six keys, it reports no tested static baseline above $0.503$. \ou\ These finite baseline fits do not imply that the all-lag representation lacks information. Sample sizes, model revisions, training budgets, exact seed identities and the definition of reported variability must be supplied before these numbers can establish a complete comparison. No exact AUC is invented for cells described only as ``at chance.''

\textbf{Additional recommendation details.} The report selects 30Music using the dataset comparison of \citet{klenitskiy2024} and states that it had the strongest shuffle-related degradation among fifteen evaluated datasets. \ou\ The report's relative SASRec and GRU4Rec losses are $86\%$ and $90\%$, respectively; we keep the underlying absolute scores in the main text because uncertainty and denominators matter. \ou\ The candidate universe, sampled-negative protocol, preprocessing filters, user/session boundary handling, target construction, split sizes and whether every perturbed condition retrains the model are not specified sufficiently in the supplied report. The claimed incomplete BERT recommender contributes no result. Recovery of those settings is R7, not a license to assume a standard protocol.

\section{Required verification and optional extensions}
\label{app:required}
These items are unfinished work, not additional completed experiments. Costs are planning estimates for one H200, following the author's corrected hardware budget. CPU-only recovery needs no GPU allocation. Future training must use separately named, corrected datasets/configurations; it must not overwrite the original results. The available budget is $300$ H200-hours, with $60$ reserved for overruns. \prov{TODO-REQUIRED}\ Recovering the reported fine-tuning, KIP and recommendation archives takes priority over new training. Rerun costs remain conditional on measured throughput and on which outputs can be recovered.

\textbf{R1: Data and implementation integrity.} \req{R1}{recover main generator versions, split hashes, all seed manifests and checkpoint versions. Independently recompute every deterministic label, the noisy flip mask and split disjointness; repair checkpoint cloning and real-token/padding separation in a new experiment version. Budget 2--6 H200-hours for timing and inference checks, plus CPU audit. Exact agreement validates the task description; unresolved mismatch prevents the affected mechanism claim. If corrected reference training is needed, use Transformer/LSTM/hybrid on the three original tasks over five seeds, with a separately declared fixed schedule (10--16 hours), and BERT on the two OC tasks (38--44 hours).}

\textbf{R2: Model-level feature diagnostics.} \req{R2}{construct 10,000 unique opposite-label OC pairs with identical letter and lag-pair histograms, using multiple key triples/residue templates and non-key fillers. Assert every label and feature equality; reserve templates before looking at neural scores. Evaluate frozen checkpoints without retuning. Estimate within-pair ranking, with half credit for ties, and bootstrap seeds/template families with 10,000 resamples. Budget 1--4 H200-hours. Above-tie performance refutes a pair-count-only account; failures disprove exact rule recovery on the diagnostic inputs.}

\textbf{R3: Scope the parity training result.} \req{R3}{recover original curriculum runs. For the planned corrected comparison, train Transformer/LSTM on raw GPC using five seeds, 100 epochs, no early stopping; compare direct length-20 training with 50 epochs at length ten followed by 50 at twenty, keeping optimizer steps matched and recording the unequal token totals. Tune learning rates $\{10^{-4},3\times10^{-4},10^{-3}\}$ in matched validation-only pilots; use constant-rate Adam, batch size $64$, and no optimizer reset between curriculum stages. Select one rate per architecture from mean final-length validation loss across both schedules, then freeze it for both arms. Use test error at most $0.01$ as the prespecified success endpoint and retain the final-epoch result. Budget 30--50 H200-hours. Any success refutes a blanket failure claim; no success supports only failure within the specified schedule. Show every seed and exact success-rate intervals.}

\textbf{R4: Check adaptation rather than infer architecture.} \req{R4}{recover full-fine-tuning and LoRA configurations, exact trainable parameter counts, model/tokenizer versions, selected checkpoints and per-sequence scores. Check split, prompt, loss, readout, learning-rate and schedule matching, including LM padding masks. Budget 0--4 H200-hours if checkpoints exist. Verified improvement supports adaptation sensitivity; mismatches require a narrower statement. A new full-fine-tuning cost is unknown until throughput/memory are measured. A separate small-Transformer attention-mask/readout factorial, two tasks and five seeds per cell, is estimated at 8--14 hours if an architecture claim is retained.}

\textbf{R5: Verify the positive KIP demonstration.} \req{R5}{recover the complete four-/six-key generator and run matrix, seed IDs, sample counts, exact ranking/label convention, checkpoints and ordered/shuffled predictions. Assert each key occurs once, independence of position assignment, disjoint splits and shuffle preservation of original labels. Recompute every reported AUC and define each variability statistic. Budget 0--4 H200-hours from archived checkpoints. A reproduced LSTM success with both shuffle controls supports order-dependent prediction on this task; failed verification suspends the neural claim, without affecting the mathematical construction. Retraining is unbudgeted pending a pilot.}

\textbf{R6: Repair the clinical comparison.} \req{R6}{inside credentialed access, keep the original patients/split, truncate once to the same retained code list, then shuffle; repair packed recurrent states. Use Transformer/LSTM/BiLSTM, five training seeds and three fixed shuffle draws, with paired initialization/minibatch order and validation-only model/threshold selection. Report AUC, average precision and paired bootstrap intervals over patients, seeds and draws, keeping predictions within the secure environment. Budget 4--12 H200-hours. A proposed practical AUC margin of 0.01 must be fixed before new results; only intervals contained within that margin support bounded small effects. Wide intervals or positive material effects refute the proposed conclusion or leave it unresolved.}

\textbf{R7: Verify the recommendation protocol.} \req{R7}{recover the six table rows' predictions and code, candidate sets, negative sampling, split/preprocessing parameters, training-versus-test shuffle behavior and exact seed/draw pairing. Check fixed visible items before comparing perturbations, and group bootstrap resampling at user level with shared predictions across conditions. Budget 0--4 H200-hours from archived models. Reproducing the fixed-content drop supports use of order under that protocol; recent-context controls test whether long histories are needed. No training cost is assumed without the model/configuration archive.}

\textbf{R8: Final reporting.} \req{R8}{complete missing run-level statistics, source/version manifests and the author-reviewed AI disclosure; regenerate all new tables from predictions, preserve failures and explain discrepancies. Use Holm correction for prespecified superiority families. Report uncertainty over seeds separately from conditional test uncertainty; five independent paired seeds cannot give a two-sided exact sign-test p-value below 0.0625. Budget 1--4 H200-hours for final scoring, otherwise CPU work. Successful deterministic regeneration verifies displayed numbers; unresolved missing data remain flagged rather than replaced by plausible estimates.}

\opt{O1}{add five confirmation seeds to the four direct/curriculum model cells, with the selected schedules unchanged, if measured savings remain; estimated 25--40 H200-hours. This improves success-rate precision but does not establish a universal learnability claim.}

\opt{O2}{repeat the rule audit across five independently selected key rules with LSTM/Transformer and three seeds, preserving the original parameters; estimated 10--20 H200-hours. This tests breadth across rules rather than zero-shot identification of an unseen label map.}

\opt{O3}{complete the missing two Llama-1B objective/readout cells, five seeds and matched validation tuning, if component attribution is needed; estimated 75--90 H200-hours. A gain in a combined recipe cannot otherwise be assigned to its objective or readout separately.}

\end{document}
